% !TeX root = ../papers/nc-p3.tex
\section{任意の中心有限正則窓の因子類群と有理二次曲面での分岐障害}
\label{sec:finite-center-divisor-torsion}
\begin{theorem}\label{thm:finite-center-divisor-torsion}
任意の成分上で標準正則窓の折り畳みが中心 $Z$ 上有限とし、頂点指数を $s$ とする。
アフィン中心錐のWeil因子類群 $\operatorname{Cl}(Z)$ は有限で $16s$ 倍で消える。
より正確に $s\operatorname{Cl}(Z)$ は $(\mathbb Z/2)^2\oplus\mathbb Z/16$ の
ある商群の部分群である。射影中心 $X$ の有理因子類群は豊富な偏極の張る一次元空間で、
$X$ は $\mathbb Q$-factorialである。有限四辺・指数四・四次中心パラメータをこの定理に仮定しない。
\end{theorem}
\begin{proof}
先の中心有限窓の定理から $Z$ は正規整域で、滑らかな部分 $U\subset\operatorname{Spec} Z$ で
折り畳みはAzumayaである。これは一般のパラメータ次数でも中心双対・有限大域次元の議論による。
$\operatorname{Cl}(Z)=\operatorname{Pic}(U)$。分裂gerbe上で四冪等元に対応する同じ階数 $s$ の
四束 $E_i$ を取り、$d_i$ をそれぞれの行列式とする。すべての差 $d_i-d_j$ は $U$ 上の通常の直線束である。
頂点単純のKoszul分解を $U$ に制限し、森田同値で束に移す。
錐の頂点を除いてよく、余次元4なのでPicard群は変わらない。ここでは経路次数の捩りを忘れているから
$E_{i+4}=E_i$ で、完全列の行列式は
\[
      d_{i+4}-4d_{i+3}+6d_{i+2}-4d_{i+1}+d_i=0.
\]
従って差の張る群は $\mathbb Z[C_4]$ の増大イデアル $I=(z-1)$ の $I/I^4$ の商である。
$u=z-1$ とすれば $u,u^2,u^3$ の関係は $4u+6u^2+4u^3=4u^2+6u^3=4u^3=0$。
Smith標準形は $2,2,16$ である。

$\operatorname{Cl}(Z)$ は斉次高さ一素因子で生成される。実際斉次非零元全部を逆にすると
次数零の体上のLaurent一変数環になる。非零次数群の正の生成元の次数の元を一つ選べばよい。
この局所化は一意分解なので、斉次元の主因子の成分が類群を生成する。各成分は次数作用で保たれ斉次である。
そこで任意の斉次反射的階数一加群 $L$ に対し、$Be_0$ と $L$ の中心上の積の二重双対を取る。
これは有限生成次数付き左 $B$ 加群で、$U$ での森田像は（側の規約を揃えて）$E_0\otimes L$ である。
有限斉次射影分解の項は移動した頂点射影の和である。分解での各頂点の符号付き個数を $n_i$ とすると、
$U$ 上で階数を比べ $\sum n_i=1$、行列式を比べ $s[L]=\sum n_i(d_i-d_0)$ を得る。
二重双対の次数は経路次数に合わせて $L$ の圧縮次数を四倍すればよい。これで指数の主張が従う。
$U$ は滑らかな有限型多様体で $H^1_{\rm et}(U,\mu_{16s})$ は有限であり、
Kummer列でPicard群のこの捩れ部分へ全射する。従って類群そのものも有限である。

各射影素因子の斉次錐について、ある正倍は斉次元の主因子である。
斉次性は主イデアルの次数作用またはその斉次生成元で分かる。底に十分高い共通次数 $N$ を取り
$\mathcal O_X(N)$ を豊富で可逆かつ生成されるものとする。上の元の適当な正冪は
$\mathcal O_X(N)$ の冪の切断となって元の素因子のみに消え、他の素因子を持たない。
従ってその素因子の正倍は豊富な偏極の正倍と線型同値である。これで射影側の結論を得る。
\end{proof}

\begin{proposition}\label{prop:quadric-pole-residue-test}
定理\ref{thm:quartic-center-image-atlas}の第二型のうち四次系に基点があり、
$Q=(q_2=0)\subset\mathbb P^4$ が滑らかな場合を取る。
$T=(l_2=0)\subset Q$ が滑らかなら、この中心を実現する正則窓は存在しない。
さらに $T$ が既約正規なら、$T$ の滑らかな部分には非自明な位数2または4の巡回指標が必要である。
\end{proposition}
\begin{proof}
射影中心は $Q$ と双有理で $l_2\ne0$ 上では $w=-f_4/l_2$ によって同型となる。
ここは滑らかで次数1の局所単元があり、零番Azumaya隅の類
$\alpha\in\operatorname{Br}(k(Q))[4]$ はここで非分岐である。その指数は4であり非零である。
滑らかな有理射影二次曲面三次元体のBrauer群は零なので、高さ一の剰余の少なくとも一つが非零で、
それは $T$ の成分でしか起きない。$T$ が既約正規なら正則三次元体 $Q$ の剰余複体で
隣り合う剰余の合成は零であるから、その唯一の剰余は $T$ の滑らかな部分の全ての因子で非分岐となる。
純性によりそこに延びる非零巡回指標である。滑らかな $T$ は豊富な超曲面で連結、
Lefschetzの基本群定理で $Q$ と同じ自明な幾何的基本群を持つので矛盾する。
\end{proof}
